\documentclass[11pt]{article}
\usepackage[margin=1in]{geometry}
\usepackage{microtype}
\usepackage{amsmath,amssymb,amsthm,mathtools}
\usepackage[hidelinks]{hyperref}
\usepackage[capitalize,nameinlink]{cleveref}

\newtheorem{definition}{Definition}[section]
\newtheorem{proposition}[definition]{Proposition}
\newtheorem{lemma}[definition]{Lemma}
\newtheorem{theorem}[definition]{Theorem}
\newtheorem{corollary}[definition]{Corollary}
\theoremstyle{remark}
\newtheorem{remark}[definition]{Remark}

\newcommand{\RR}{\mathbb{R}}
\newcommand{\ZZ}{\mathbb{Z}}
\newcommand{\CP}{\mathcal{C}}
\newcommand{\eps}{\varepsilon}
\newcommand{\norm}[1]{\left\lVert #1 \right\rVert}
\newcommand{\abs}[1]{\left\lvert #1 \right\rvert}
\newcommand{\set}[1]{\left\{ #1 \right\}}
\newcommand{\odd}{\operatorname*{\sum\nolimits_{i\ \mathrm{odd}}}}
\newcommand{\even}{\operatorname*{\sum\nolimits_{i\ \mathrm{even}}}}
\newcommand{\Pass}{\mathrm{Pass}}
\newcommand{\Fail}{\mathrm{Fail}}
\newcommand{\Warn}{\mathrm{Warn}}
\newcommand{\NotRun}{\mathrm{NotRun}}
\newcommand{\Unknown}{\mathrm{Unknown}}

\title{CP Generator: Formal Notes}
\author{Allan Pan}
\date{\today}

\begin{document}

\maketitle

\section{The implemented mathematical object}

The project works with a square sheet
\[
S=[0,s]^2,\qquad s>0,
\]
and with a finite straight-line graph drawn on that sheet. More precisely, the
implemented state is a triple
\[
\CP=(G,X,\tau),
\]
where $G=(V,E)$ is a finite graph, $X\colon V\to S$ is an embedding, and
\[
\tau\colon E\to\set{-1,0,1}
\]
is a fold-label map, with $-1$ for ``unassigned'', $0$ for ``mountain'', and
$1$ for ``valley''.

\begin{definition}[Boundary and interior vertices]
\[
V_{\partial}=\set{v\in V:X_v\in\partial S},
\qquad
V_{\mathrm{int}}=V\setminus V_{\partial}.
\]
\end{definition}

Every local theorem used by the program is expressed at interior vertices, so
the distinction between $V_{\partial}$ and $V_{\mathrm{int}}$ is structural
rather than cosmetic.

\begin{definition}[Local angle data]
For $v\in V_{\mathrm{int}}$, let
\[
d(v)=\deg_G(v).
\]
If the incident creases at $v$ are listed in cyclic order, then the
corresponding sector angles are written
\[
\alpha_1(v),\dots,\alpha_{2m(v)}(v)\in(0,2\pi),
\qquad
\sum_{i=1}^{2m(v)}\alpha_i(v)=2\pi.
\]
\end{definition}

\begin{definition}[Kawasaki residual and Maekawa balance]
For $v\in V_{\mathrm{int}}$, define
\[
\kappa(v)
=
\max\!\left(
\abs{\odd \alpha_i(v)-\pi},
\abs{\even \alpha_i(v)-\pi}
\right).
\]
If $\tau$ is complete at $v$, define
\[
M_\tau(v)=\#\set{e\ni v:\tau(e)=0},
\qquad
V_\tau(v)=\#\set{e\ni v:\tau(e)=1},
\qquad
\mu_\tau(v)=M_\tau(v)-V_\tau(v).
\]
\end{definition}

\section{From random points to a planar candidate graph}

The first stage of the project is purely geometric. One samples interior points
\[
P\subset (0,s)^2,
\]
adjoins the four corners of the square, and then projects points lying within a
fixed threshold $\delta>0$ of the boundary onto $\partial S$. Thus one obtains
\[
P^\square=P\cup\set{(0,0),(s,0),(s,s),(0,s)},
\qquad
\widetilde P=\Pi_{\partial,\delta}(P^\square).
\]

The candidate graph is then taken from the Delaunay triangulation of
$\widetilde P$:
\[
G_0=\operatorname{Del}(\widetilde P).
\]
After this, the code removes every edge lying entirely on one side of the
square boundary. Therefore the implemented crease graph is
\[
E=E(G_0)\setminus E_{\partial},
\]
where
\[
E_{\partial}
=
\set{
\set{u,v}\in E(G_0):
X_u,X_v \text{ lie on one boundary side of }S
}.
\]

This Delaunay stage is not itself a flat-foldability theorem. Its role is to
provide a planar graph with reasonably regular local geometry so that the later
angle constraints are numerically meaningful.

\section{Workflow-level genericity guard}

The local flat-foldability equations alone do not prevent a numerically
collapsed sample in which two vertices become arbitrarily close. The workflow
therefore adds a scale-normalized geometric genericity test before it accepts a
sample as ``green''.

\begin{definition}[Minimum vertex separation]
For an embedded pattern $\CP=(G,X,\tau)$ on a square of side length $s$, define
\[
\delta_{\min}(X)
=
\min_{\substack{u,v\in V\\u\ne v}}
\norm{X_u-X_v},
\qquad
\rho(X)=\frac{\delta_{\min}(X)}{s}.
\]
\end{definition}

\begin{definition}[Implemented generic-geometry predicate]
The workflow fixes the ratio threshold
\[
\rho_{\min}=10^{-3}.
\]
It calls the geometry \emph{generic} when
\[
\delta_{\min}(X)\ge \rho_{\min}s.
\]
Equivalently,
\[
\operatorname{Generic}(X)
\Longleftrightarrow
\rho(X)\ge 10^{-3}.
\]
\end{definition}

The pattern generator retries up to a fixed maximum number of random samples
and returns the first generic one. If no generic sample is found, it keeps the
best available sample, namely the one maximizing $\delta_{\min}(X)$ among those
attempts. Likewise, the local optimizer is declared ``green'' only when
\[
\operatorname{Local}(\CP;\eps)
\quad\text{and}\quad
\operatorname{Generic}(X)
\]
hold simultaneously.

\begin{definition}[Workflow-level local green]
The workflow-level local success predicate is
\[
\operatorname{LocalGreen}(\CP;\eps)
\Longleftrightarrow
\operatorname{Local}(\CP;\eps)\wedge \operatorname{Generic}(X).
\]
\end{definition}

\begin{remark}
This genericity guard is not a theorem from classical origami theory. It is an
implemented numerical nondegeneracy assumption, used to reject nearly
coincident vertices for which angle diagnostics and later exact-face reasoning
become unstable.
\end{remark}

\section{The local flat-foldability conditions}

The next stage turns the geometric candidate graph into a constrained origami
object by imposing the standard single-vertex necessary conditions.

\begin{theorem}[Single-vertex necessary conditions {\cite{hull_math_flat}}]
If $v\in V_{\mathrm{int}}$ is flat-foldable as a single-vertex crease pattern,
then
\[
d(v)\in 2\ZZ,
\qquad
\odd \alpha_i(v)=\pi,
\qquad
\even \alpha_i(v)=\pi,
\qquad
\abs{\mu_\tau(v)}=2.
\]
\end{theorem}

The optimizer in the code imposes only one alternating-angle equality per
vertex, not two. This is sufficient, because the second equality is redundant.

\begin{lemma}[Alternating-sum redundancy]
If
\[
\sum_{i=1}^{2m}\alpha_i=2\pi,
\]
then
\[
\odd \alpha_i=\pi
\Longleftrightarrow
\even \alpha_i=\pi.
\]
\end{lemma}

\begin{proof}
Since
\[
\even \alpha_i = 2\pi-\odd \alpha_i,
\]
each equality implies the other.
\end{proof}

Accordingly, the project defines its local diagnostic predicate by
\[
\operatorname{Local}(\CP;\eps)
\Longleftrightarrow
\forall v\in V_{\mathrm{int}}\;
\bigl(
d(v)\in 2\ZZ
\ \wedge\
\kappa(v)\le \eps
\bigr).
\]
This is the mathematical content of the program's \emph{local status}.

\section{Parity repair before optimization}

The first obstruction to local flat-foldability is odd degree. Let
\[
O(G)=\set{v\in V:d(v)\notin 2\ZZ}
\]
denote the odd-degree set. The implementation repairs parity by repeatedly
deleting shortest paths between odd vertices.

\begin{definition}[Implemented parity-repair step]
Choose
\[
u\in O(G),
\qquad
w=\operatorname*{arg\,min}_{z\in O(G)\setminus\set{u}} d_G(u,z),
\]
and let
\[
\gamma=(u=v_0,v_1,\dots,v_k=w)
\]
be a shortest path from $u$ to $w$ in $G$. Replace $E$ by
\[
E\setminus \set{\set{v_{i-1},v_i}:1\le i\le k}.
\]
\end{definition}

The reason this step is mathematically natural is that path deletion affects
parity only at the endpoints.

\begin{proposition}[Parity effect of one deletion]
If the path $\gamma$ above is deleted, then
\[
\deg(v_i)\mapsto \deg(v_i)-2
\qquad (1\le i\le k-1),
\]
and
\[
\deg(u)\mapsto \deg(u)-1,
\qquad
\deg(w)\mapsto \deg(w)-1.
\]
Hence the parity of every internal path vertex is preserved, while the parity
of the endpoints flips.
\end{proposition}

\begin{proof}
Each internal path vertex loses exactly two incident edges, and each endpoint
loses exactly one.
\end{proof}

This proves that the repair has the intended local parity effect. It does not
prove that the repaired graph is globally optimal in any stronger sense.

\section{Weighted geometric repair via constrained optimization}

Once parity has been repaired, the code moves the vertices as little as possible
while imposing Kawasaki's identity at every interior vertex.

Let $X^{(0)}=(X_v^{(0)})_{v\in V}$ denote the pre-optimization coordinates.
Boundary motion is heavily penalized by weights
\[
w_v=
\begin{cases}
10^6,& v\in V_{\partial},\\
1,& v\in V_{\mathrm{int}}.
\end{cases}
\]
The weighted displacement functional is
\[
D(X)=\sum_{v\in V} w_v \norm{X_v-X_v^{(0)}}^2.
\]

The actual code minimizes
\[
F(X)=\frac{0.1}{2}D(X)^2.
\]
Since $D(X)\ge 0$, one has
\[
\arg\min F = \arg\min D.
\]
Thus the implementation still solves the nearest-feasible-geometry problem in
the weighted least-squares sense.

For each interior vertex $v$, define the constraint
\[
g_v(X)=\odd \alpha_i(v;X)-\pi.
\]
The implemented optimization problem is therefore
\begin{equation}\label{eq:opt}
\min_X F(X)
\qquad
\text{subject to}
\qquad
g_v(X)=0
\quad
\forall v\in V_{\mathrm{int}}.
\end{equation}

The solver is SLSQP. After the numerical solve, vertices within a small
tolerance of the boundary are snapped back onto $\partial S$. In this way, the
project passes from a merely planar sample to a planar sample constrained by
the local angle theorem.

\section{Mountain--valley search by Bern--Hayes-style local reduction}

After geometric repair, the remaining local problem is to assign mountain and
valley labels. The code does not search over all edge labelings directly.
Instead, it first extracts local sign relations from smallest sectors, following
the Bern--Hayes reduction idea {\cite{bern_hayes}}.

For an interior vertex $v$, let
\[
e_1,\dots,e_r
\]
be the incident creases in cyclic order, and let
\[
\beta_1,\dots,\beta_r
\]
be the corresponding sector angles. The reduction repeatedly removes a locally
minimal sector and records whether its two boundary creases must carry the same
or opposite mountain--valley sign.

\begin{definition}[Local reduction chain]
The reduction at $v$ produces
\[
\Gamma_v=\bigl((f_1,\dots,f_k),(\eta_1,\dots,\eta_k)\bigr),
\qquad
\eta_i\in\set{0,1,2},
\]
where
\[
\eta_i=
\begin{cases}
0,& \tau(f_i)=1-\tau(f_{i-1}),\\
1,& \tau(f_i)=\tau(f_{i-1}),\\
2,& \tau(f_i)\text{ remains undetermined}.
\end{cases}
\]
\end{definition}

Boundary vertices are treated similarly, except that temporary virtual boundary
rays are added so that the cyclic order closes. Thus the local reduction stage
produces one or more chains of sign relations at every vertex.

\section{Exact enumeration after chain merging}

The local chains are then merged across shared creases. If the merged chains are
\[
\Lambda_1,\dots,\Lambda_g,
\]
then each $\Lambda_j$ contributes one free binary choice, namely the sign of
its first crease. Therefore the remaining search space is
\[
\set{0,1}^g.
\]

For
\[
c=(c_1,\dots,c_g)\in\set{0,1}^g,
\]
the code propagates the implied signs along each chain:
\[
\tau_c(e_{j,1})=c_j,
\]
and then
\[
\eta_{j,i}=0 \Longrightarrow \tau_c(e_{j,i})=1-\tau_c(e_{j,i-1}),
\]
\[
\eta_{j,i}=1 \Longrightarrow \tau_c(e_{j,i})=\tau_c(e_{j,i-1}),
\]
\[
\eta_{j,i}=2 \Longrightarrow \tau_c(e_{j,i})=-1.
\]

This yields a partial or complete labeling $\tau_c$. The exact admissibility
filter is Maekawa's theorem:
\[
\operatorname{Adm}(c)
\Longleftrightarrow
\forall v\in V_{\mathrm{int}}\;
\abs{\mu_{\tau_c}(v)}=2.
\]

Among admissible choices, the code chooses the one maximizing an interlacing
score
\[
\operatorname{Score}(c)
=
0.05\,\#\set{e\in E:\tau_c(e)\in\set{0,1}}
\sum_{v\in V}\Phi_v(c)-4\,G(\tau_c),
\]
where $\Phi_v(c)$ rewards alternating mountain--valley patterns more than
repeating ones, with sharper wedges weighted more heavily, and where
$G(\tau_c)$ counts the obtuse-angle defects defined in the next section. This
score is not a correctness condition; it is only a tie-breaker among already
admissible assignments.

\begin{proposition}[What the assignment search certifies]
If the routine returns an assignment $\tau^\ast$, then
\[
\forall v\in V_{\mathrm{int}}\;
\abs{\mu_{\tau^\ast}(v)}=2.
\]
If the routine returns failure, then no enumerated merged-group assignment
satisfies Maekawa's theorem at every interior vertex.
\end{proposition}

\begin{proof}
The code checks every seed choice in $\set{0,1}^g$, rejects every choice
violating Maekawa, and returns only from the surviving set.
\end{proof}

\section{Maekawa-preserving repair of obtuse monochrome glitches}

Local Maekawa admissibility is still not enough to rule out every visibly bad
mountain--valley pattern. The implementation now singles out one specific local
configuration that repeatedly produced globally implausible flat stacks.

\begin{definition}[Obtuse monochrome glitch]
Let $v\in V_{\mathrm{int}}$, and let
\[
e_i,e_{i+1},e_{i+2}
\]
be three consecutive incident creases in cyclic order, with consecutive sector
angles
\[
\beta_i,\beta_{i+1}.
\]
For a complete local assignment $\tau$, this triple is called an
\emph{obtuse monochrome glitch} if
\[
\beta_i>\frac{\pi}{2},
\qquad
\beta_{i+1}>\frac{\pi}{2},
\qquad
\tau(e_i)=\tau(e_{i+2})\neq \tau(e_{i+1}).
\]
\end{definition}

Equivalently, two consecutive obtuse wedges are bounded by a
``same--different--same'' sign pattern. In the user-facing language of the
program, the three creases should instead become monochromatic.

\begin{definition}[Two monochromizing moves]
For a glitch triple $(e_i,e_{i+1},e_{i+2})$, the code tests exactly two local
updates:
\[
\text{(middle flip)}\qquad
\tau'(e_{i+1})=\tau(e_i),
\]
and
\[
\text{(outer flip)}\qquad
\tau''(e_i)=\tau''(e_{i+2})=\tau(e_{i+1}),
\]
with every other crease unchanged.
\end{definition}

Only one of these moves is accepted, and only when it preserves Maekawa's
theorem not merely at $v$, but at every interior vertex.

\begin{definition}[Accepted glitch repair]
For a glitch triple $(e_i,e_{i+1},e_{i+2})$ at $v$, an update
\[
\tau\mapsto\widetilde\tau
\]
is \emph{accepted} if it is one of the two monochromizing moves above and
\[
\abs{\mu_{\widetilde\tau}(w)}=2
\qquad
\forall w\in V_{\mathrm{int}}.
\]
\end{definition}

The repair loop then chooses among all accepted updates by the lexicographic
rule
\[
\bigl(G(\widetilde\tau),\ \#\text{changed creases},\ \operatorname{Sig}(\widetilde\tau)\bigr),
\]
where smaller glitch count wins first, then fewer changed creases, then the
assignment signature used by the code for deterministic tie-breaking. The loop
repeats until no accepted update decreases $G$ or a previously seen signature
would be revisited.

\begin{proposition}[What the repair loop guarantees]
If the obtuse-glitch repair routine terminates with assignment $\widehat\tau$,
then every accepted intermediate update preserved
Maekawa's theorem at all interior vertices, and the final assignment has no
smaller reachable glitch count among the one-step accepted updates considered
from the final state.
\end{proposition}

\begin{proof}
Every candidate update is explicitly checked against Maekawa's theorem on every
interior vertex before acceptance. The loop only commits an update when it
strictly improves the current glitch count, and it halts once no such improving
accepted update exists or a repeated signature would occur.
\end{proof}

\section{From local certificates to global geometry}

At this point, the project has local angle control and a locally admissible
mountain--valley labeling. The next question is global: whether the current
straight-line embedding can be organized into coherent folded faces.

The first global obstruction is crossing.

\begin{definition}[Crossing predicate]
For distinct nonincident edges $e,e'\in E$, define
\[
\operatorname{Cross}(e,e';X)
\Longleftrightarrow
X(e)\cap X(e')\ne\varnothing.
\]
Set
\[
\operatorname{Noncross}(X)
\Longleftrightarrow
\forall e\ne e' \in E,\quad
e\cap e'=\varnothing
\Longrightarrow
\neg \operatorname{Cross}(e,e';X).
\]
\end{definition}

If $\operatorname{Noncross}(X)$ fails, then no exact face-based folded preview
is attempted on the current geometry.

\section{Exact face extraction by half-edge tracing}

Assume now that the current embedding is noncrossing. The code augments the fold
graph by the four boundary edges of the square:
\[
E^\square = E \cup E_{\partial S}.
\]
At each vertex $v$, neighbors are ordered by polar angle around $X_v$. This
induces a left-face successor rule on oriented edges.

\begin{definition}[Left-face successor]
For an oriented edge $(u,v)$, let
\[
\sigma(u,v)=(v,w),
\]
where $w$ is the predecessor of $u$ in the cyclic order around $v$.
\end{definition}

Starting from any half-edge and iterating $\sigma$ yields a closed walk or a
failure. Every bounded counterclockwise closed walk is declared to be a face.
Repeated vertices or nonclosing walks cause the exact solver to abort.

Thus the code produces a finite face family
\[
\mathcal F.
\]
For each fold edge $e\in E$, one then computes its face-incidence set
\[
\operatorname{Inc}(e)=\set{f\in\mathcal F:e\subset\partial f}.
\]
The exact solver requires
\[
\abs{\operatorname{Inc}(e)}=2
\qquad
\forall e\in E.
\]
This is the precise global topological condition needed before reflection
propagation can begin.

\section{Reflection propagation and exact-current-geometry certification}

Suppose every fold edge borders exactly two faces. The code chooses a root face
\[
f_0\in\mathcal F,
\qquad
\operatorname{area}(f_0)=\max_{f\in\mathcal F}\operatorname{area}(f),
\]
and sets
\[
T_{f_0}=I_3.
\]

If faces $f$ and $g$ share a fold line $\ell$, then the child transform is
defined by reflection:
\[
T_g = R_\ell T_f.
\]
When a face is reached by two different reflection paths, the resulting
transforms are compared on the vertices of that face.

\begin{definition}[Cycle discrepancy]
For a face $f$ and two candidate transforms $T,T'$, define
\[
\Delta(f;T,T')
=
\max_{v\in \operatorname{Vert}(f)}
\norm{T(X_v)-T'(X_v)}.
\]
\end{definition}

If
\[
\Delta(f;T,T') > \Delta_{\mathrm{hard}},
\]
then exact reconstruction fails. If
\[
\Delta_{\mathrm{soft}} < \Delta(f;T,T') \le \Delta_{\mathrm{hard}},
\]
then the reconstruction survives only at warning level.

Fold signs are resolved from $\tau$ by
\[
\tau(e)=0 \Longrightarrow \operatorname{sgn}(e)=+1,
\qquad
\tau(e)=1 \Longrightarrow \operatorname{sgn}(e)=-1.
\]
If a fold remains unassigned, the exact solver fills its sign provisionally by a
face-depth rule; this also downgrades the output to warning level.

\begin{theorem}[Exact-current-geometry certificate]
If the project reports
\[
\operatorname{Status}_{\mathrm{global}}(\CP)=\Pass
\qquad\text{and}\qquad
\operatorname{Status}_{\mathrm{preview}}(\CP)=\Pass,
\]
then
\[
\operatorname{Noncross}(X),
\qquad
\abs{\operatorname{Inc}(e)}=2 \ \forall e\in E,
\qquad
\Delta(f;T,T')\le \Delta_{\mathrm{soft}}
\]
for every cycle comparison, and no provisional fold signs were used.
\end{theorem}

\begin{proof}
Each clause is the negation of one of the exact solver's failure or warning
branches. A pass can occur only if every one of those branches is avoided.
\end{proof}

This theorem describes the strongest global statement currently certified by the
code. It is a theorem about the \emph{current embedding}, not a complete
classification theorem for all embeddings with the same combinatorics.

\section{Bounded exact stack-order checking}

The exact preview solver reconstructs faces and rigid transforms, but it does
not by itself enumerate all possible flat overlap orders. For sufficiently
small exact instances, the project therefore runs a separate bounded exact
checker on the extracted face arrangement.

\begin{definition}[Bounded-check prerequisites]
The bounded checker runs only if all of the following hold:
\begin{enumerate}
    \item $\operatorname{Status}_{\mathrm{local}}(\CP)=\Pass$,
    \item $\operatorname{Status}_{\mathrm{assign}}(\CP)=\Pass$,
    \item $\operatorname{Noncross}(X)$,
    \item exact face reconstruction succeeds,
    \item no provisional fold signs are used,
    \item no approximate cycle closure is used, and
    \item the extracted face count satisfies
    \[
    \abs{\mathcal F}\le F_{\max},
    \qquad
    F_{\max}=8
    \]
    in the current implementation.
\end{enumerate}
\end{definition}

The checker compares faces in a near-flat, but still slightly separated, rigid
pose. The implemented angle is
\[
\theta_{\mathrm{nf}}=\pi-0.08.
\]
Let
\[
P_f^{\mathrm{nf}}\subset\RR^3
\]
be the points of face $f$ at angle $\theta_{\mathrm{nf}}$, and let
\[
\overline P_f\subset\RR^2
\]
be the corresponding exact final flat projection.

\begin{definition}[Overlap cell]
Let $T\subset \overline P_f$ and $U\subset \overline P_g$ be triangles from the
face triangulations. Their \emph{overlap cell} is the convex polygon
\[
C(T,U)=T\cap U
\]
computed by polygon clipping. Cells of area at most
\[
\eps_{\mathrm{area}}=10^{-8}
\]
are discarded.
\end{definition}

For each nontrivial overlap cell, the code samples the centroid
\[
c(T,U)\in C(T,U)
\]
and interpolates the two near-flat heights
\[
z_f\bigl(c(T,U)\bigr),
\qquad
z_g\bigl(c(T,U)\bigr)
\]
from the corresponding 3-dimensional triangles by barycentric coordinates.

\begin{definition}[Pairwise overlap-order relation]
For two distinct faces $f,g$, if every overlap cell between them satisfies
\[
z_f(c)-z_g(c)>\eps_z,
\qquad
\eps_z=10^{-7},
\]
then the checker records the relation
\[
f\succ g.
\]
If every overlap cell instead satisfies
\[
z_g(c)-z_f(c)>\eps_z,
\]
it records
\[
g\succ f.
\]
If some overlap cells reverse the sign, or if
\[
\abs{z_f(c)-z_g(c)}\le \eps_z
\]
on a nontrivial overlap, the checker fails immediately.
\end{definition}

Thus the checker constructs a directed relation
\[
R\subset \mathcal F\times\mathcal F,
\]
interpreted as a partial overlap order on faces. The bounded exact problem is
then whether $R$ admits a linear extension.

\begin{definition}[Bounded stack order]
A \emph{bounded stack order} for the extracted arrangement is a permutation
\[
(f_{i_1},\dots,f_{i_m})
\]
of $\mathcal F$ such that
\[
(f,g)\in R
\Longrightarrow
f \text{ appears before } g.
\]
\end{definition}

The implementation counts bounded stack orders by a bitmask dynamic program. If
\[
\operatorname{Pred}(f)=\set{g:(g,f)\in R},
\]
and if $N(S)$ denotes the number of admissible completions from the already
placed face set $S\subseteq\mathcal F$, then
\[
N(S)=
\sum_{\substack{f\in \mathcal F\setminus S\\ \operatorname{Pred}(f)\subseteq S}}
N(S\cup\set{f}),
\qquad
N(\mathcal F)=1.
\]
The count is truncated above by
\[
N_{\max}=100000.
\]

\begin{proposition}[Meaning of a bounded exact-check pass]
If the bounded checker returns $\Pass$, then the extracted relation $R$ admits
at least one bounded stack order. If it returns $\Pass$ with unique-order flag,
then exactly one bounded stack order was found before the truncation limit was
reached.
\end{proposition}

\begin{proof}
The dynamic program above counts precisely the linear extensions of the
predecessor relation encoded by $R$, truncated only when the count exceeds
$N_{\max}$. A positive count therefore certifies existence of at least one
linear extension, and count $1$ certifies uniqueness.
\end{proof}

\begin{proposition}[Meaning of a cyclic bounded failure]
If the bounded checker returns failure with message ``cyclic stack-order
constraint'', then the extracted relation $R$ has no linear extension, hence no
bounded stack order compatible with all sampled overlap constraints.
\end{proposition}

\begin{proof}
The cited failure branch is reached exactly when the topological-order counter
returns $0$. A directed acyclic predecessor relation has at least one linear
extension, so zero count means that the extracted overlap-order constraints are
incompatible.
\end{proof}

\begin{remark}
This bounded checker is still not a complete global flat-foldability algorithm.
It is a small-instance certificate on the exact face arrangement actually
constructed by the solver, and it is philosophically closest to the
overlap-order viewpoint used in the flat-folding complexity literature
{\cite{bern_hayes,eppstein_flat}}.
\end{remark}

\section{3D folding and fallback preview}

Once exact face reconstruction is available, the project builds a kinematic 3D
folding by propagating rotations along the face-adjacency tree. If a crease $e$
has axis $(p,q)$ and sign $\operatorname{sgn}(e)\in\set{\pm1}$, then the
3-dimensional transform satisfies
\[
T_g^{(3D)}
=
R_{(p,q)}\!\bigl(\operatorname{sgn}(e)\theta\bigr)T_f^{(3D)}.
\]
This rigid-hinge viewpoint is the standard rigid-plate abstraction used for the
preview stage; compare the rigid-origami literature
{\cite{abel_rigid,tachi_simulation}}.

For time parameter $t\in[0,1]$, the implementation uses the easing profile
\[
\eta(t)=\frac{1-\cos(\pi t)}{2},
\]
and for any settle window $[a,b]\subseteq[0,1]$ it uses the smoothstep blend
\[
s_{a,b}(t)=
\begin{cases}
0,& t\le a,\\[2mm]
3u(t)^2-2u(t)^3,& a\le t\le b,\quad
u(t)=\dfrac{t-a}{b-a},\\[2mm]
1,& t\ge b.
\end{cases}
\]
Let $\Pi:\RR^3\to\RR^2$ be projection to the $(x,y)$-plane, let $X_f$ denote
the ordered source vertices of face $f$ with initial $z$-coordinate $0$, and
define the two exact rigid guides
\[
P_f^{\mathrm{tree}}(t)
=
T_f^{(3D)}\bigl((\pi-0.1)\eta(t)\bigr)(X_f),
\qquad
P_f^{\mathrm{exact}}(t)
=
T_f^{(3D)}\bigl(\pi\eta(t)\bigr)(X_f).
\]
We also write
\[
P^{\mathrm{tree}}(t)=\bigl(P_f^{\mathrm{tree}}(t)\bigr)_{f\in\mathcal F},
\qquad
P^{\mathrm{exact}}(t)=\bigl(P_f^{\mathrm{exact}}(t)\bigr)_{f\in\mathcal F}
\]
for the corresponding facewise families.
The first is the capped tree pose used for the display-oriented legacy and
balanced modes; the second is the full exact folded guide used by the seamless
rigid mode.

\begin{definition}[Rigid face fit]
For a face $f$ and a target point set
\[
Y_f=(y_{f,1},\dots,y_{f,m_f})\in(\RR^3)^{m_f},
\]
the implemented rigid fit is
\[
\operatorname{RigidFit}(X_f,Y_f)=R_fX_f+u_f,
\]
where $(R_f,u_f)\in SO(3)\times\RR^3$ minimizes
\[
\sum_{i=1}^{m_f}\norm{R_fx_{f,i}+u_f-y_{f,i}}^2.
\]
Numerically this is the usual orthogonal Procrustes / Kabsch solve with the
reflection branch removed.
\end{definition}

\begin{lemma}[Rigid fit preserves face geometry]
For every $f$ and every target point set $Y_f$, the point set
\[
\operatorname{RigidFit}(X_f,Y_f)
\]
is congruent to $X_f$. In particular, for every pair of indices $i,j$,
\[
\norm{\bigl(\operatorname{RigidFit}(X_f,Y_f)\bigr)_i
-\bigl(\operatorname{RigidFit}(X_f,Y_f)\bigr)_j}
=
\norm{x_{f,i}-x_{f,j}}.
\]
\end{lemma}

\begin{proof}
By definition,
\[
\bigl(\operatorname{RigidFit}(X_f,Y_f)\bigr)_i
=
R_fx_{f,i}+u_f
\]
with $R_f\in SO(3)$. Hence
\[
\norm{(R_fx_{f,i}+u_f)-(R_fx_{f,j}+u_f)}
=
\norm{R_f(x_{f,i}-x_{f,j})}
=
\norm{x_{f,i}-x_{f,j}},
\]
because orthogonal matrices preserve Euclidean norm.
\end{proof}

\subsection*{Legacy layered}

\paragraph{Intuition.}
Legacy layered is the most conservative display mode. For most of the animation
it simply shows the original tree-based rigid fold. Only near the end does it
peel the faces apart into a visibly separated flat stack, so that the final
layer order is easy to read even when many faces would otherwise collapse onto
the same plane. The price is that the final peel-away is a display blend, not a
physically rigid motion.

\begin{definition}[Legacy layered stack]
Let $m=\abs{\mathcal F}$. Evaluate the nearly-flat exact pose at angle
\[
\theta_{\mathrm{nf}}=\pi-0.08.
\]
For each face $f$, let $(x_f^{\mathrm{nf}},y_f^{\mathrm{nf}},z_f^{\mathrm{nf}})$
be the centroid of
\[
T_f^{(3D)}(\theta_{\mathrm{nf}})(X_f).
\]
Order the
faces by the lexicographic key
\[
\bigl(z_f^{\mathrm{nf}},y_f^{\mathrm{nf}},x_f^{\mathrm{nf}}\bigr),
\]
and let $r(f)\in\set{0,\dots,m-1}$ be the resulting rank. With
\[
\operatorname{span}(X)
=
\bigl(\max_i X_i^{(1)}-\min_i X_i^{(1)}\bigr)
+
\bigl(\max_i X_i^{(2)}-\min_i X_i^{(2)}\bigr),
\]
define the implemented legacy spacing
\[
\Delta_{\mathrm{leg}}
=
\max\left\{
\frac{\operatorname{span}(X)}{18m},
0.02
\right\},
\qquad
h_f^{\mathrm{leg}}
=
\Delta_{\mathrm{leg}}
\left(r(f)-\frac{m-1}{2}\right).
\]
If $\Phi_f^{\mathrm{flat}}:\RR^2\to\RR^2$ is the exact flat transform of face
$f$, the legacy target stack is
\[
L_f
=
\Bigl\{
\bigl(\Phi_f^{\mathrm{flat}}(\Pi(x_{f,i})),h_f^{\mathrm{leg}}\bigr)
\Bigr\}_{i=1}^{m_f}.
\]
The rendered legacy motion is
\[
P_f^{\mathrm{legacy}}(t)
=
\bigl(1-s_{0.72,1}(t)\bigr)P_f^{\mathrm{tree}}(t)
+
s_{0.72,1}(t)L_f.
\]
\end{definition}

\begin{proposition}[Legacy layered endpoints]
For every face $f$,
\[
P_f^{\mathrm{legacy}}(0)=X_f,
\qquad
P_f^{\mathrm{legacy}}(t)=P_f^{\mathrm{tree}}(t)
\text{ for }0\le t\le 0.72,
\qquad
P_f^{\mathrm{legacy}}(1)=L_f.
\]
Moreover the path $t\mapsto P_f^{\mathrm{legacy}}(t)$ is continuous on
$[0,1]$.
\end{proposition}

\begin{proof}
Since $\eta(0)=0$, one has
\[
P_f^{\mathrm{tree}}(0)=T_f^{(3D)}(0)(X_f)=X_f.
\]
By definition of $s_{0.72,1}$, the blend factor vanishes on $[0,0.72]$ and
equals $1$ at $t=1$. Therefore the displayed endpoint identities follow
immediately. Continuity follows from continuity of $\eta$, $s_{0.72,1}$,
$P_f^{\mathrm{tree}}$, and affine interpolation in Euclidean space.
\end{proof}

\begin{proposition}[Legacy layered is generally not rigid during the settle window]
Fix $t\in(0.72,1)$ and abbreviate $s=s_{0.72,1}(t)\in(0,1)$. For two vertices
$a,b$ of the same face, let
\[
\Delta_{ab}^{\mathrm{tree}}
=
P_{f,a}^{\mathrm{tree}}(t)-P_{f,b}^{\mathrm{tree}}(t),
\qquad
\Delta_{ab}^{\mathrm{flat}}
=
L_{f,a}-L_{f,b},
\]
and let $\ell_{ab}=\norm{x_{f,a}-x_{f,b}}$. Then
\[
\norm{P_{f,a}^{\mathrm{legacy}}(t)-P_{f,b}^{\mathrm{legacy}}(t)}^2
=
\ell_{ab}^2
+
2s(1-s)\bigl(
\Delta_{ab}^{\mathrm{tree}}\cdot\Delta_{ab}^{\mathrm{flat}}
-\ell_{ab}^2
\bigr).
\]
In particular, unless the endpoint edge vectors coincide for every edge of the
face, the late legacy blend is not a rigid motion.
\end{proposition}

\begin{proof}
Both $P_f^{\mathrm{tree}}(t)$ and $L_f$ are rigid images of $X_f$, hence
\[
\norm{\Delta_{ab}^{\mathrm{tree}}}
=
\norm{\Delta_{ab}^{\mathrm{flat}}}
=
\ell_{ab}.
\]
The blended edge vector is
\[
(1-s)\Delta_{ab}^{\mathrm{tree}}+s\Delta_{ab}^{\mathrm{flat}},
\]
so expanding its squared norm gives
\begin{align*}
\norm{(1-s)\Delta_{ab}^{\mathrm{tree}}+s\Delta_{ab}^{\mathrm{flat}}}^2
&=
(1-s)^2\ell_{ab}^2+s^2\ell_{ab}^2
+2s(1-s)\Delta_{ab}^{\mathrm{tree}}\cdot\Delta_{ab}^{\mathrm{flat}}\\
&=
\ell_{ab}^2
+2s(1-s)\bigl(
\Delta_{ab}^{\mathrm{tree}}\cdot\Delta_{ab}^{\mathrm{flat}}
-\ell_{ab}^2
\bigr).
\end{align*}
This equals $\ell_{ab}^2$ for all $s\in(0,1)$ only in the exceptional case that
every edge vector agrees between the two endpoint placements.
\end{proof}

\subsection*{Balanced stack}

\paragraph{Intuition.}
Balanced stack starts from the same tree pose, but it tries to heal visible
cracks before the final display blend. The program repeatedly averages the
different copies of each shared vertex and then rigidly refits each whole face
to those averaged targets. This does not attempt to recover the exact physical
folded state. Its purpose is more modest and more visual: reduce the worst seam
openings while preserving a small amount of display thickness in the final
stack.

For a global vertex $v$, let
\[
\mathcal O(v)=
\set{(f,i):\text{the $i$th local vertex of face $f$ represents }v}.
\]
Given a facewise point collection $Q=(Q_f)_f$, define the averaged shared
vertex location
\[
a_v(Q)
=
\frac{1}{\abs{\mathcal O(v)}}
\sum_{(f,i)\in\mathcal O(v)}Q_{f,i}.
\]
For $0\le\lambda\le1$, one relaxation sweep forms targets
\[
\widetilde Q_{f,i}
=
(1-\lambda)Q_{f,i}+\lambda a_{v_{f,i}}(Q)
\]
and then refits the whole face by
\[
\bigl(\mathcal U_\lambda(Q)\bigr)_f
=
\operatorname{RigidFit}(X_f,\widetilde Q_f).
\]
The $k$-sweep relaxation operator is
\[
\mathcal R_{k,\lambda}
=
\mathcal U_\lambda^k.
\]

\begin{definition}[Balanced stack target and motion]
Using the same rank order $r(f)$ as legacy layered, the thinner balanced stack
uses spacing
\[
\Delta_{\mathrm{bal}}
=
\max\left\{
\frac{\operatorname{span}(X)}{220m},
0.006
\right\},
\qquad
h_f^{\mathrm{bal}}
=
\Delta_{\mathrm{bal}}
\left(r(f)-\frac{m-1}{2}\right),
\]
and flat target
\[
S_f
=
\Bigl\{
\bigl(\Phi_f^{\mathrm{flat}}(\Pi(x_{f,i})),h_f^{\mathrm{bal}}\bigr)
\Bigr\}_{i=1}^{m_f}.
\]
Let
\[
C_f
=
\bigl(\mathcal R_{6,0.9}(P^{\mathrm{tree}}(1))\bigr)_f,
\qquad
B_f
=
0.64\,C_f+0.36\,S_f.
\]
At runtime the balanced mode uses
\[
Q_f^{\mathrm{bal}}(t)
=
\bigl(\mathcal R_{k(t),0.88}(P^{\mathrm{tree}}(t))\bigr)_f,
\qquad
k(t)=
\begin{cases}
4,& t<0.65,\\
5,& t\ge 0.65,
\end{cases}
\]
followed by the final display blend
\[
P_f^{\mathrm{bal}}(t)
=
\bigl(1-s_{0.76,1}(t)\bigr)Q_f^{\mathrm{bal}}(t)
+
s_{0.76,1}(t)B_f.
\]
\end{definition}

\begin{proposition}[Balanced relaxation preserves per-face rigidity]
For every point collection $Q$, every face of
\[
\mathcal R_{k,\lambda}(Q)
\]
is congruent to the source face $X_f$.
\end{proposition}

\begin{proof}
Each sweep $\mathcal U_\lambda$ replaces a face by
\[
\operatorname{RigidFit}(X_f,\widetilde Q_f),
\]
which is congruent to $X_f$ by the rigid-fit lemma. Iterating the same fact
$k$ times proves the claim.
\end{proof}

\begin{proposition}[Closed rigid configurations are fixed points of the relaxation]
Suppose that for each face $f$ there are $(R_f,u_f)\in SO(3)\times\RR^3$ with
\[
Q_f=R_fX_f+u_f,
\]
and suppose moreover that whenever two local vertices represent the same global
vertex, their points in $Q$ coincide exactly. Then for every
$\lambda\in[0,1]$,
\[
\mathcal U_\lambda(Q)=Q.
\]
\end{proposition}

\begin{proof}
If all copies of a global vertex already coincide, then
\[
a_v(Q)=Q_{f,i}
\]
for every $(f,i)\in\mathcal O(v)$, so the intermediate targets satisfy
\[
\widetilde Q_{f,i}=Q_{f,i}.
\]
Thus $\mathcal U_\lambda$ rigidly fits $X_f$ to the exact rigid image $Q_f$
itself. Because every exact face has positive area, at least three noncollinear
points determine the rigid transform uniquely, so the zero-residual fit returns
$Q_f$ unchanged.
\end{proof}

\begin{corollary}[Balanced endpoints]
For every face $f$,
\[
P_f^{\mathrm{bal}}(0)=X_f,
\qquad
P_f^{\mathrm{bal}}(1)=B_f.
\]
\end{corollary}

\begin{proof}
At $t=0$ one has $P_f^{\mathrm{tree}}(0)=X_f$, and all shared copies of each
vertex coincide. The preceding proposition therefore gives
\[
Q_f^{\mathrm{bal}}(0)=X_f.
\]
Also $s_{0.76,1}(0)=0$ and $s_{0.76,1}(1)=1$, so the displayed endpoint formula
follows immediately from the final blend definition.
\end{proof}

\begin{remark}
As in legacy layered mode, the last affine blend toward $B_f$ is generally not
itself a rigid motion; the rigid guarantee applies to the relaxation output
$Q_f^{\mathrm{bal}}(t)$ before the final display blend.
\end{remark}

\subsection*{Seamless rigid}

\paragraph{Intuition.}
Seamless rigid is the exact-only mode that tries to remove both visual defects
at once: it wants connected faces in the middle of the motion, and it wants the
last frame to be the true folded figure rather than a display stack. The exact
fold is therefore used only as a guide. At discrete sample times the code
re-solves a new family of rigid face transforms that stays close to the exact
guide, keeps neighboring hinge rotations compatible with that guide, and
explicitly penalizes seam openings. The displayed motion is then built only from
those accepted rigid samples.

\begin{definition}[Seamless rigid samples]
Let $r$ be the root face and let
\[
t_j=\frac{j}{12},
\qquad
j=0,\dots,12.
\]
For each sample time $t_j$, define the exact guide points
\[
G_f(t_j)=P_f^{\mathrm{exact}}(t_j),
\]
and let
\[
\operatorname{RigidFit}(X_f,G_f(t_j))
=
\widehat R_{f,j}X_f+\widehat u_{f,j}
\]
be the facewise guide fit. The root face is fixed at
\[
R_{r,j}=I,
\qquad
u_{r,j}=0.
\]
For each global vertex $v$, let $(f_v,i_v)$ denote the first stored occurrence
in $\mathcal O(v)$. The three residual families are
\begin{align}
E^{\mathrm{seam}}_j
&=
\sum_v
\sum_{(f,i)\in\mathcal O(v)\setminus\set{(f_v,i_v)}}
\norm{
(R_{f,j}x_{f,i}+u_{f,j})
-
(R_{f_v,j}x_{f_v,i_v}+u_{f_v,j})
}^2,
\notag\\
E^{\mathrm{guide}}_j
&=
\sum_{f\ne r}\sum_i
\norm{R_{f,j}x_{f,i}+u_{f,j}-G_{f,i}(t_j)}^2,
\notag\\
E^{\mathrm{hinge}}_j
&=
\sum_{(f,g)\in H}
\norm{
\log\!\Bigl(
(R_{g,j}R_{f,j}^{\top})
(\widehat R_{g,j}\widehat R_{f,j}^{\top})^{\top}
\Bigr)
}^2,
\notag
\end{align}
where $H$ is the set of fold-adjacent face pairs. The implemented objective is
\begin{align}
\mathcal E_j
&=
130\,E^{\mathrm{seam}}_j
+0.34\,E^{\mathrm{guide}}_j
+0.75\,E^{\mathrm{hinge}}_j.
\label{eq:seamless-rigid}
\end{align}
If $p^\star$ denotes the parameter vector of the exact final fit
$F_f^\star=P_f^{\mathrm{exact}}(1)$ and $\widehat p_j$ the guide-fit parameter
vector, then the optimizer is started from
\[
p_j^{(0)}
=
\begin{cases}
\widehat p_j,& t_j\le0.84,\\[2mm]
(1-t_j)\widehat p_j+t_jp^\star,& t_j>0.84.
\end{cases}
\]
\end{definition}

If an accepted sample has incident copies
\[
(a_e^+,b_e^+),
\qquad
(a_e^-,b_e^-)
\]
of a shared edge $e$, define the edge mismatch
\[
\operatorname{gap}_e
=
\min\Bigl\{
\max\{\norm{a_e^+-a_e^-},\norm{b_e^+-b_e^-}\},
\max\{\norm{a_e^+-b_e^-},\norm{b_e^+-a_e^-}\}
\Bigr\}.
\]
The sample is accepted exactly when
\[
\operatorname{Gap}_j
=
\max_e\operatorname{gap}_e
\le 0.02.
\]

Between two solved neighboring samples $t_a<t_b$, the rendered in-between
rotation of each face is the geodesic / SLERP interpolation
\[
R_f(t)
=
R_{f,a}
\exp\!\Bigl(
\alpha(t)\log\bigl(R_{f,a}^{\top}R_{f,b}\bigr)
\Bigr),
\qquad
\alpha(t)=\frac{t-t_a}{t_b-t_a},
\]
with linearly interpolated translation
\[
u_f(t)
=
(1-\alpha(t))u_{f,a}+\alpha(t)u_{f,b}.
\]
If one of the immediate neighboring samples failed, the renderer holds the
nearest solved exact sample instead of dropping to the older gappy panel
relaxation.

\begin{proposition}[Accepted seamless-rigid samples are facewise rigid and seam-certified]
If sample $t_j$ is accepted, then every face of that sample is congruent to
$X_f$, and its maximum shared-edge mismatch satisfies
\[
\operatorname{Gap}_j\le 0.02.
\]
\end{proposition}

\begin{proof}
Accepted samples are stored only in the form
\[
R_{f,j}X_f+u_{f,j},
\]
so every face is a rigid image of $X_f$. The bound on $\operatorname{Gap}_j$
is exactly the acceptance criterion used by the implementation.
\end{proof}

\begin{proposition}[Every interpolated seamless-rigid frame is facewise rigid]
Suppose $t\in(t_a,t_b)$ lies between two solved neighboring samples. Then for
every face $f$, the rendered in-between points have the form
\[
R_f(t)X_f+u_f(t)
\]
with $R_f(t)\in SO(3)$. Hence each face remains congruent to $X_f$ throughout
the interpolated interval.
\end{proposition}

\begin{proof}
The matrix
\[
\exp\!\Bigl(
\alpha(t)\log\bigl(R_{f,a}^{\top}R_{f,b}\bigr)
\Bigr)
\]
lies in $SO(3)$, so $R_f(t)\in SO(3)$. Therefore
\[
R_f(t)X_f+u_f(t)
\]
is a rigid image of $X_f$, and the rigid-fit lemma gives congruence of all
facewise distances.
\end{proof}

\begin{proposition}[Seamless rigid endpoints and exact-model fallback behavior]
For the exact model,
\[
P_f^{\mathrm{rigid}}(0)=X_f,
\qquad
P_f^{\mathrm{rigid}}(1)=F_f^\star=P_f^{\mathrm{exact}}(1).
\]
Moreover the implementation stores solved anchor samples at $t_0=0$ and
$t_{12}=1$ from initialization, so every progress value is rendered from either
a solved exact sample or an interpolation between solved exact samples; it does
not need to switch back to the older legacy panel relaxation.
\end{proposition}

\begin{proof}
The code explicitly stores the initial source configuration as solved sample
$t_0$ and the exact folded fit as solved sample $t_{12}$. It also returns those
two states directly at $t=0$ and $t=1$. Since these two anchors always exist,
every intermediate progress value has at least one solved sample at or before
it and at least one solved sample at or after it. Hence the exact renderer can
always return either a solved exact sample or an interpolation between solved
exact samples.
\end{proof}

\subsection*{Triangle-mesh fallback}

\paragraph{Intuition.}
If exact face reconstruction fails, the program does not claim a certified
rigid-face motion. Instead it triangulates the available regions and propagates
rigid motions along a triangle tree. This keeps a useful picture on screen, but
the resulting preview is explicitly warning-level and heuristic.

The mesh fallback triangulates each polygonal region into triangles
\[
\Delta.
\]
Let $\theta_{\mathrm{mesh}}^{\max}$ denote the mesh preview angle cap, equal to
$76^\circ$ in the flattening case and $88^\circ$ otherwise. If a tree edge
between triangles corresponds to a fold edge $e$ with approximate normalized
dihedral weight $\alpha_e$, then the child transform is
\[
S_{\Delta'}(t)
=
R_{(p,q)}\!\bigl(\theta_{\mathrm{mesh}}^{\max}\eta(t)\,\alpha_e\bigr)
S_\Delta(t).
\]
If the tree edge is merely a triangulation edge, the child simply inherits the
parent transform. The legacy and balanced display variants are then applied to
these triangle poses in exactly the same style as above.

\begin{proposition}[Triangle rigidity in mesh fallback]
Before the late display blend, every triangle in the mesh fallback remains
rigid, and every child triangle shares its tree edge exactly with its parent.
\end{proposition}

\begin{proof}
Each triangle transform is obtained from its parent either by applying a rigid
rotation about the common edge or by copying the parent transform unchanged.
Hence every triangle pose is a rigid image of its source triangle, and the
common tree edge is preserved exactly by construction.
\end{proof}

\begin{remark}
The mesh fallback remains intentionally non-certifying: it does not prove exact
global face reconstruction, exact cycle closure, or exact overlap order.
\end{remark}

\section{Status predicates}

The four displayed statuses are designed to keep the local and global
statements separate.

\begin{definition}[Implemented status semantics]
\[
\operatorname{Status}_{\mathrm{local}}(\CP)
=
\begin{cases}
\Pass,& \operatorname{Local}(\CP;\eps),\\
\Fail,& \text{otherwise},
\end{cases}
\]
\[
\operatorname{Status}_{\mathrm{assign}}(\CP)
=
\begin{cases}
\Fail,& \exists v\in V_{\mathrm{int}}:\abs{\mu_\tau(v)}\ne 2
\text{ with complete local assignment data},\\
\NotRun,& \forall e\in E:\tau(e)=-1,\\
\Warn,& \exists e\in E:\tau(e)=-1 \text{ and some assignments exist},\\
\Pass,& \operatorname{Assign}(\CP),
\end{cases}
\]
\[
\operatorname{Status}_{\mathrm{global}}(\CP)
=
\begin{cases}
\Fail,& \neg\operatorname{Noncross}(X),\\
\NotRun,& E=\varnothing,\\
\Unknown,& \operatorname{Status}_{\mathrm{assign}}(\CP)=\NotRun,\\
\Pass,\Warn,\Fail,& \text{exact-current-geometry solver output},
\end{cases}
\]
\[
\operatorname{Status}_{\mathrm{preview}}(\CP)
=
\begin{cases}
\NotRun,& \text{no preview attempt admissible},\\
\Pass,\Warn,\Fail,& \text{preview solver output}.
\end{cases}
\]
\end{definition}

The automation layer also uses two compound predicates.

\begin{definition}[Workflow-level green predicates]
The local optimization loop uses $\operatorname{LocalGreen}$ above, while the
full automation predicate is
\[
\begin{aligned}
\operatorname{AllGreen}(\CP)
\Longleftrightarrow\;&
\operatorname{Status}_{\mathrm{local}}(\CP)=\Pass
\wedge
\operatorname{Status}_{\mathrm{assign}}(\CP)=\Pass
\\
&\wedge
\operatorname{Status}_{\mathrm{global}}(\CP)=\Pass
\wedge
\operatorname{Status}_{\mathrm{preview}}(\CP)=\Pass.
\end{aligned}
\]
\end{definition}

\begin{remark}
Thus the word ``green'' has two different implemented meanings: one local and
geometric, used during optimization; and one global four-status conjunction,
used by the full search workflows.
\end{remark}

\section{Scope and limitations}

The project combines rigorous local constraints with a partly rigorous and
partly heuristic global pipeline. The following distinctions are therefore
essential.

\begin{itemize}
    \item The shortest-path parity repair is an implemented heuristic, not an
    optimality theorem.
    \item The Bern--Hayes-style wedge reduction is a local sign-reduction
    device derived from {\cite{bern_hayes}}, not a complete multi-vertex
    flat-foldability classifier.
    \item The exact face-reflection solver certifies the current embedding when
    it passes, but the mesh fallback is only a guarded visualization.
    \item Warning-level outputs indicate exactly where the code leaves the realm
    of strict certification.
\end{itemize}

\begin{theorem}[Implemented scope]
The current pipeline is not a complete decision procedure for arbitrary
multi-vertex global flat-foldability.
\end{theorem}

\begin{proof}
The claim follows from the coexistence of heuristic parity repair, local-only
sign reduction, and explicit mesh/reference fallbacks in the preview stage.
\end{proof}

\begin{thebibliography}{99}

\bibitem{bern_hayes}
Marshall Bern and Barry Hayes.
\newblock \emph{The Complexity of Flat Origami}.
\newblock SODA, 1996.

\bibitem{hull_math_flat}
Thomas C. Hull.
\newblock \emph{On the Mathematics of Flat Origamis}.
\newblock Congressus Numerantium, 1994.

\bibitem{abel_rigid}
Zachary Abel, Jason Cantarella, Erik D. Demaine, David Eppstein, Thomas C.
Hull, Jason S. Ku, Robert J. Lang, and Tomohiro Tachi.
\newblock \emph{Rigid Origami Vertices: Conditions and Forcing Sets}.
\newblock Journal of Computational Geometry, 2017.

\bibitem{tachi_simulation}
Tomohiro Tachi.
\newblock \emph{Simulation of Rigid Origami}.
\newblock In \emph{Origami$^4$}, 2009.

\bibitem{eppstein_flat}
David Eppstein.
\newblock \emph{A Parameterized Algorithm for Flat Folding}.
\newblock arXiv:2306.11939, 2023.

\bibitem{hull_zakharevich}
Thomas C. Hull and Inna Zakharevich.
\newblock \emph{Flat Origami Is Turing Complete}.
\newblock arXiv:2309.07932, version consulted in 2025.

\bibitem{walker_stankovic}
Andreas Walker and Tino Stankovi\'c.
\newblock \emph{Algorithmic Design of Origami Mechanisms and Tessellations}.
\newblock Communications Materials, 2022.

\end{thebibliography}

\end{document}
